\section{Implementation and Investigation Notes}
\label{sec:implementation}

\subsection{Activation Supply and Cache Lifetime}
Training consumes local saved trajectories, with activations constructed on
demand by an external vLLM extraction server. A bounded CPU cache retains full
tensors, while prompt summaries can be reused separately. Cache capacity is
not a cap on process RSS: temporary construction buffers, model state, and
other allocations contribute additional memory. Persistent summaries,
prefetching, and concurrency controls are implementation features; only the
pilot timings in SYS-001 are treated as measurements here.

\subsection{Checkpoint Publication and Recovery}
The full-data training workflow supports local recovery and Hub snapshots
containing generator weights, optimizer state, progress, RNG state, resolved
configuration, and available provenance. Inference bundles can be downloaded
without optimizer checkpoints. Recovery capability is recorded as workflow
support; it does not demonstrate completion of the ten-epoch training plan.

\subsection{Independent Serving and Benchmarking}
The suite separates model serving from the GuideLLM client. It launches and
cleans up owned server/client process groups, exports category workloads once,
records attempt state, and can skip completed measurements with matching
artifact checksums. Model type selects an original drafter, static-LoRA bundle,
or generator bundle. Throughput can be measured without an additional
correctness-generation pass; the absence of that pass is recorded explicitly.

\subsection{Fixed-Batch Integration Boundaries}
The new scheduler leaves vacant slots empty until a group finishes. It permits
a partial group after a 50 ms collection window in the example configuration.
The runtime rejects KV preemption or incomplete initial group admission rather
than reuse draft cache entries under a different generated adapter. Larger
concurrency also increases the prefill/profile token budget and may change
GPU memory requirements. These constraints define BATCH-001
and the shared-adapter BATCH-002 protocol. GPU summary measurements are now
recorded, while detailed scheduling traces remain unavailable.

Collection, queueing, uneven request completion, and partial tails are part of
the measured serving protocol. Closed groups are not equivalent to unrestricted
continuous batching. Batch membership may vary with arrival/completion order,
even when the same prompt list is reused across model variants.

\subsection{Per-Request Multi-LoRA and Native Baseline}
The per-request controller generates factors once from complete target prefill,
retains a request-specific CPU copy, and activates only scheduled request IDs in
the draft-only native LoRA manager. Token mappings follow actual query offsets
on every draft forward, including speculative rejection and later one-token
steps. Request completion releases that request's adapter; paused request factors
remain available. The native scheduler can refill free slots without changing
other requests' adapters. CPU tests cover lifecycle and mapping invariants;
aggregate BATCH-003 counters are not a full isolation or correctness test.

Baseline launches can create an isolated package view containing the original,
checksum-verified copies of the three integration-patched vLLM files. The global
installation is left intact, and worker status can report imported source hashes.
This supports original-source EAGLE3 with native admission. Such source hashes
were not included in the supplied BATCH-003 upload, so this capability must not
be presented as independently verified provenance for those measurements.

\subsection{Bounded Profiling: Implemented, Measurements Pending}
Revision \texttt{dc872c5} adds a two-phase profiling command. An unprofiled pass
provides ordinary timing; fresh servers then collect bounded native CPU/CUDA
traces after warmup on the same frozen models and workload. The serial example
compares original EAGLE3, singleton-conditioned merged adapters and singleton
native multi-LoRA; the continuous example compares baseline and per-request
adapters at C4. Example captures use math/chat, 16 prompts per category,
128-token caps and at most 30 worker iterations.

Custom ranges cover generation, packing/loading, mapping and draft/projection
execution. Analysis records CPU inclusive/self time, GPU kernels and transfers,
union activity and gaps, and captured request occupancy. Projection ranges
include base GEMM work; isolated LoRA kernels require their own kernel rows.
The native execute-model scope includes target work and preparation rather than
providing an isolated verification timer. Unattributed activity remains unknown.
A bounded trace may miss refills or completions, and whole-workload acceptance
counters cannot supply captured-window accepted-token denominators.

CPU lifecycle tests, synthetic CUDA-event analysis, a real CPU profiler export
and mocked HTTP/GuideLLM orchestration checks were performed locally. CUDA/vLLM
profiling has not been executed locally and no user GPU traces have been supplied.
Profiler timings are diagnostic and excluded from ordinary throughput tables.
Thus recurring adapter kernels/routing remain hypotheses, not diagnosed causes.
